\section*{Цель работы}
Исследование температурных зависимостей емкости, тангенса угла диэлектрических потерь и температурного коэффициента диэлектрической проницаемости линейных и нелинейных диэлектриков.

\section*{Теоретические сведения}
Диэлектрики --- это материалы или среды с удельным сопротивлением более $10^8 \text{ Ом}\cdot\text{м}$, в которых при приложении электрического поля возникает эффект поляризации. Поляризация может быть вызвана упругим смещением электронных оболочек (электронная), ориентацией диполей (дипольно-релаксационная), смещением ионов (ионная) или упорядочением доменов (спонтанная). 

Относительная диэлектрическая проницаемость характеризует способность диэлектриков поляризоваться:
\begin{equation}
    \varepsilon = \frac{C_{\text{д}}}{C_0},
\end{equation}
где $C_{\text{д}}$ --- емкость конденсатора с диэлектриком; $C_0$ --- емкость того же конденсатора в вакууме.

При работе конденсатора в переменном электрическом поле часть энергии рассеивается в виде тепла. Полные потери в участке изоляции с емкостью $C$ при напряжении $U$ и круговой частоте $\omega$ определяются выражением:
\begin{equation}
    P_a = U^2 \omega C \tan \delta,
\end{equation}
где $\delta$ --- угол диэлектрических потерь (угол, дополняющий до $90^\circ$ угол сдвига фаз между током и напряжением). Параметр $\tan \delta$ характеризует способность диэлектрика рассеивать энергию.

Емкость плоского конденсатора зависит от геометрических размеров и свойств диэлектрика:
\begin{equation}
    C = \frac{\varepsilon \varepsilon_0 S}{h},
\end{equation}
где $\varepsilon_0 = 8.85 \cdot 10^{-12} \text{ Ф/м}$ --- электрическая постоянная; $S$ --- площадь электродов; $h$ --- толщина диэлектрика.

Температурный коэффициент емкости $\alpha_C$ отражает относительное изменение емкости при изменении температуры:
\begin{equation}
    \alpha_C = \frac{1}{C} \frac{dC}{dT}.
    \label{form:alphac}
\end{equation}

Для металлизированных твердотельных конденсаторов, где тонкий слой металла нанесен непосредственно на диэлектрик, изменение размеров электродов определяется линейным расширением самого диэлектрика. В этом случае температурный коэффициент емкости связан с температурным коэффициентом диэлектрической проницаемости $\alpha_\varepsilon$ следующим образом:
\begin{equation}
    \alpha_C = \alpha_\varepsilon + \alpha_{l\text{д}},
    \label{form:alphae}
\end{equation}
где $\alpha_{l\text{д}}$ --- температурный коэффициент линейного расширения диэлектрика.

\section*{Характеристика и области применения исследуемых материалов}

В данной работе рассматриваются диэлектрики, различающиеся по химическому составу, структуре и механизмам поляризации, что определяет их эксплуатационные характеристики.

\begin{enumerate}
    \item \textit{Неорганическое стекло}.
    Представляет собой аморфный диэлектрик с неплотной упаковкой ионов.
    Основные области применения: изготовление изоляторов, подложек интегральных схем и корпусов полупроводниковых приборов.
    В конденсаторостроении стекло используется для создания низкочастотных конденсаторов с относительно высокими потерями, где не требуется высокая температурная стабильность.
    Химический состав: аморфные оксиды.
    $\varepsilon \approx 6.5$.
    $\tan\delta = 0.0011$.
    $\alpha_\varepsilon = 0.80 \cdot 10^{-4} \dots 5.05 \cdot 10^{-4}\,\text{К}^{-1}$.

    \item \textit{Слюда}.
    Природный слоистый диэлектрик, обладающий высокой электрической прочностью и низкими потерями.
    Применяется в производстве высокочастотных и высоковольтных конденсаторов, а также в качестве межвитковой изоляции в электрических машинах большой мощности.
    Благодаря высокой нагревостойкости (до $500 \div 600\,^\circ\text{C}$) слюда незаменима в устройствах, работающих в жестких температурных условиях.
    Химический состав: природный алюмосиликат.
    $\varepsilon \approx 7.5$.
    $\tan\delta = 0.0014$.
    $\alpha_\varepsilon = -0.13 \cdot 10^{-4} \div 3.91 \cdot 10^{-4}\,\text{К}^{-1}$.

    \item \textit{Тиконд}.
    Керамический материал на основе диоксида титана ($TiO_2$, рутил).
    Применяется для изготовления высокочастотных конденсаторов, которые позволяют стабилизировать резонансную частоту колебательных контуров при изменении температуры.
    Химический состав: $TiO_2$.
    $\varepsilon \approx 55$.
    $\tan\delta = 0.0007$.
    $\alpha_\varepsilon = -15.54 \cdot 10^{-4} \div -5.86 \cdot 10^{-4}\,\text{К}^{-1}$.

    \item \textit{Полипропилен}.
    Неполярный полимерный диэлектрик с предельно низким значением тангенса угла потерь.
    Используется в производстве пленочных конденсаторов для силовой электроники, импульсных цепей и прецизионных измерительных устройств.
    Обладает высокой стабильностью параметров в широком диапазоне частот, однако ограничен по температуре эксплуатации (обычно до $100 \dots 105\,^\circ\text{C}$).
    Химический состав: $(C_3H_6)_n$.
    $\varepsilon \approx 2.2$.
    $\tan\delta = 0.0005$.
    $\alpha_\varepsilon = -10.81 \cdot 10^{-4} \div -1.71 \cdot 10^{-4}\,\text{К}^{-1}$.

    \item \textit{Сегнетокерамика}.
    Сложный оксидный материал (часто на основе титаната бария $BaTiO_3$), обладающий спонтанной поляризацией.
    Благодаря высоким значениям $\varepsilon$ (до $10^4$ и выше) сегнетокерамика применяется для создания малогабаритных конденсаторов большой емкости.
    Также используется в нелинейных устройствах, датчиках давления и пьезоэлектрических преобразователях.
    Существенным ограничением является сильная зависимость параметров от температуры и напряженности электрического поля.
    Химический состав: на основе $BaTiO_3$.
    $\varepsilon_{max} \approx 420$.
    $\tan\delta = 0.0102$.
    $\alpha_\varepsilon = -556.60 \cdot 10^{-4} \div 1005.25 \cdot 10^{-4}\,\text{К}^{-1}$.
\end{enumerate}

\section*{Описание установки}
Испытательная установка состоит из пульта и цифрового измерителя иммитанса (емкости и $\tan \delta$). В испытательном модуле расположен термостат, температура в котором задается регулятором \textit{«Установка температуры»} и контролируется термопарой, подключенной к измерительному прибору. В термостате размещены пять конденсаторов с различными рабочими диэлектриками: неорганическое стекло, слюда, тиконд, полипропилен, сегнетокерамика. Выводы конденсаторов подключены к переключателю на лицевой панели.

\section*{Порядок выполнения работы}

\begin{enumerate}
    \item \textit{Подготовка к испытаниям.}
    \begin{itemize}
        \item Включите цифровой измеритель емкости и $\tan \delta$, подготовьте его к работе.
        \item Измерьте начальную емкость соединительных проводов $C_0$ (без подключения образцов). Запишите это значение в соответствующее поле протокола.
    \end{itemize}
    
    \item \textit{Измерение параметров при комнатной температуре.}
    \begin{itemize}
        \item Убедитесь, что регулятор \textit{«Установка температуры»} находится в крайнем левом положении.
        \item Включите на пульте тумблер \textit{«Сеть»} и зафиксируйте комнатную температуру по встроенному прибору. Занесите значение в первую строку \cref{tab:measurements}.
        \item Устанавливая переключатель образцов последовательно в положения от \textit{«1»} до \textit{«5»}, измерьте емкость $C$ и $\tan \delta$ для каждого материала. 
        \item \textit{Нюанс:} Истинная емкость образца определяется как разность показаний прибора и емкости проводов $C_0$. Результаты заносите в \cref{tab:measurements}.
    \end{itemize}
    
    \item \textit{Снятие температурных зависимостей.}
    \begin{itemize}
        \item Включите на пульте тумблер \textit{«Нагрев»}.
        \item Переведите регулятор \textit{«Установка температуры»} в следующее положение (на один шаг вправо).
        \item Дождитесь стабилизации температуры (через $2\dots3$ мин после прекращения роста показаний термометра).
        \item Зафиксируйте установившуюся температуру в \cref{tab:measurements}.
        \item Поочередно переключая каналы от \textit{«1»} до \textit{«5»}, измерьте и запишите значения $C$ и $\tan \delta$ для всех образцов при данной температуре.
        \item Повторяйте ступенчатое увеличение температуры вплоть до крайнего правого положения регулятора, каждый раз дожидаясь стабилизации и снимая показания для всех пяти конденсаторов.
    \end{itemize}
    
    \item \textit{Завершение работы.}
    \begin{itemize}
        \item Верните регулятор температуры в крайнее левое положение.
        \item Выключите измеритель емкости, тумблеры \textit{«Нагрев»} и \textit{«Сеть»}.
    \end{itemize}
\end{enumerate}

\section*{Указания по обработке результатов}
\begin{enumerate}
    \item По экспериментальным данным из \cref{tab:measurements} постройте графики температурных зависимостей емкости $C(t)$ и тангенса угла диэлектрических потерь $\tan \delta(t)$ для всех пяти образцов.
    \item Вычислите значения температурного коэффициента емкости $\alpha_C$. Производную $dC/dt$ следует находить путем графического дифференцирования построенных кривых $C(t)$. Для линейных зависимостей достаточно $4\dots5$ точек, для нелинейных --- $8\dots10$ точек (обязательно включая экстремумы и точки перегиба).
    \item Рассчитайте значения температурного коэффициента диэлектрической проницаемости $\alpha_\varepsilon$ для исследованных материалов, используя формулу $\alpha_\varepsilon = \alpha_C - \alpha_{l\text{д}}$. Справочные значения $\alpha_{l\text{д}}$ берутся из методических указаний.
    \item Постройте температурные зависимости $\alpha_\varepsilon(t)$ и на их основе сделайте выводы о преобладающих механизмах поляризации в каждом из исследованных диэлектриков.
\end{enumerate}